\documentclass[12pt,reqno]{article}

\usepackage[usenames]{color}
\usepackage{amssymb}
\usepackage{amsmath}
\usepackage{amsthm}
\usepackage{amsfonts}
\usepackage{amscd}
\usepackage{graphicx}

\usepackage[colorlinks=true,
linkcolor=webgreen,
filecolor=webbrown,
citecolor=webgreen]{hyperref}

\definecolor{webgreen}{rgb}{0,.5,0}
\definecolor{webbrown}{rgb}{.6,0,0}

\usepackage{color}
\usepackage{fullpage}
\usepackage{float}

\usepackage{psfig}
\usepackage{graphics}
\usepackage{latexsym}
\usepackage{epsf}
\usepackage{breakurl}

\setlength{\textwidth}{6.5in}
\setlength{\oddsidemargin}{.1in}
\setlength{\evensidemargin}{.1in}
\setlength{\topmargin}{-.1in}
\setlength{\textheight}{8.4in}

\newcommand{\seqnum}[1]{\href{https://oeis.org/#1}{\rm \underline{#1}}}

\usepackage{mathtools}
\usepackage{booktabs}

\newcommand{\abs}[1]{\ensuremath{\left|{#1}\right|}}
\newcommand\floor[1]{\left\lfloor#1\right\rfloor}
\newcommand\pbrace[1]{\left\{#1\right\}}
\newcommand{\paren}[1]{\left\lparen{#1}\right\rparen}
\newcommand{\C}{\ensuremath{\mathbb{C}}}
\newcommand{\R}{\ensuremath{\mathbb{R}}}
\newcommand{\N}{\ensuremath{\mathbb{N}}}
\newcommand{\Z}{\ensuremath{\mathbb{Z}}}
\newcommand{\vv}[1]{\ensuremath{\mathbf{#1}}}
\newcommand{\norm}[1]{\ensuremath{\left\|{#1}\right\|}}
\newcommand{\set}[1]{\ensuremath{\left\{{#1}\right\}}}

\newlength{\minuslength}
\settowidth{\minuslength}{$-$}
\newcommand{\hm}{\ensuremath{\hspace{\minuslength}}}

\def\modd#1 #2{#1\ {\rm (mod}\ #2{\rm )}}

\begin{document}

\begin{center}
\epsfxsize=4in
\leavevmode\epsffile{logo129.eps}
\end{center}

\theoremstyle{plain}
\newtheorem{theorem}{Theorem}
\newtheorem{corollary}[theorem]{Corollary}
\newtheorem{lemma}[theorem]{Lemma}
\newtheorem{proposition}[theorem]{Proposition}

\theoremstyle{definition}
\newtheorem{definition}[theorem]{Definition}
\newtheorem{example}[theorem]{Example}
\newtheorem{conjecture}[theorem]{Conjecture}

\theoremstyle{remark}
\newtheorem{remark}[theorem]{Remark}

\begin{center}
\vskip 1cm{\LARGE\bf 
Unital Sums of the M\"{o}bius \\
\vskip .1in
and Mertens Functions
}
\vskip 1cm
\large
Jeffery Kline\\
Madison, WI 53711\\
USA\\
\href{mailto:jeffery.kline@gmail.com}{\tt jeffery.kline@gmail.com}
\end{center}

\vskip .2 in

\begin{abstract}
We use standard techniques of linear algebra to construct an infinite 
family of identities that involve finite weighted sums of the 
M\"{o}bius and Mertens functions, where the weights are equal 
to $-1$, $0$, or $1$.
In a related manner, we construct, for each positive integer $n$, 
an $n\times n$ symmetric unimodular matrix, and each matrix is used to 
express an identity that involves a finite weighted sum of 
the M\"{o}bius function.  We establish several results on the spectral 
decomposition of these matrices.
\end{abstract}

\section{Introduction}
\label{sec:intro}
The M\"{o}bius function is defined as
\begin{align*}
\mu(k)\coloneqq \begin{cases}
\phantom{-}1, & \text{if $k$ is squarefree and has an even number of prime factors;} \\
-1,           & \text{if $k$ is squarefree and has an odd number of prime factors;} \\
\phantom{-}0, & \text{otherwise.}
\end{cases}
\end{align*}
Summatory behavior of $\mu$ is of great interest, in part due to the close connection that $\mu$ has to the
distribution of the prime numbers. Indeed, it is known that the Mertens function, which we write as $M(n)\coloneqq\sum_{k\leq n}\mu(k)$, 
satisfies $M(n)=o(n)$, and that this asymptotic bound is equivalent to the
prime number theorem. 

It is the goal of the first part of this paper to state and prove several summatory identities satisfied by $\mu$ and $M$ that
we have not found in prior work. These identities are stated in Proposition~\ref{prp:main} and
Corollary~\ref{cor:main} of Section~\ref{sec:main}. In Section~\ref{sec:qn}, we employ similar techniques to construct, for 
each positive integer $n$, an $n\times n$ symmetric matrix, $Q_n$, having integer entries and determinant 1 that 
satisfies $\sum_{k\leq n} Q_{n}(j,k)\mu(k) = 0$ when $j>1$, and the sum equals~1 when $j=1$.
This is the content of Proposition~\ref{prp:qn}.  We also obtain asymptotic bounds on the extremal eigenvalues of $Q_n$.
This is the content of Proposition~\ref{prp:lambdanbound}. A corollary of these propositions
is the asymptotic bound $\sum_{i\leq n} v_n(i)\mu(i)\leq n^{-3/2}C(1+o(1))$, where 
$v_n$ is the normalized dominant eigenvector of $Q_n$, and $C$ is a constant smaller than 3.  
This is stated in Corollary~\ref{cor:mun}.  We conclude with a conjecture about $v_n$.

We now briefly summarize Proposition~\ref{prp:main} and Corollary~\ref{cor:main} of the next section. For each $m,n\geq 1$, we construct $(-1,0,1)$-valued functions, 
$\epsilon^{m,n}$ and $E^{m,n}$, that satisfy the identities
\begin{align*}
\sum_{k\leq n} \epsilon^{m,n}(k)\mu(k)=1
\end{align*}
and
\begin{align*}
\sum_{k\leq n} E^{m,n}(k)M(k)=1.
\end{align*}
Each function $\epsilon^{m,n}$ is the result of applying the transpose of Dirichlet convolution by the $n$-term 
sequence $(1,1,1,\ldots,1)$ to a particular $4m$-periodic function.
Each function $E^{m,n}$ is a finite difference operation applied to $\epsilon^{m,n}$. Complete details are presented in Section~\ref{sec:main}.
Before proceeding, we state a concise formula for 
$\epsilon^{m,n}$ in the special case $m=1$. For all $k\geq 1$,
\begin{align}
\epsilon^{1,n}(k) = \sum_{ j k\leq n} \cos\paren{\frac{\pi (jk-1)}{2}}.
\label{eqn:cos}
\end{align}
See Section~\ref{subsec:cos} for further observations concerning the function $\epsilon^{1,n}$.

\begin{figure*}[ht]
\begin{center}
\includegraphics[width=0.85\textwidth]{imgexmaple}
\end{center}
\caption{Examples of the function $\epsilon^{m,n}$.}
\label{fig:main}
\end{figure*}
Figure~\ref{fig:main} displays $\epsilon^{m,n}$ for various choices of $m$ and~$n$. 
Both $E^{m,n}$ and $\epsilon^{m,n}$ are $4m$-periodic on the interval $n/3<k\leq n$.
The behavior of $E^{m,n}$ and $\epsilon^{m,n}$ near the origin is 
more complicated.   In general, $\epsilon^{m,n}$ and $E^{m,n}$ are sparse when $m$ is large.

The body of literature on the functions $\mu$ and $M$ is vast, 
and many results on sums with the form 
\begin{align*}\sum_{k\leq n}w(k)\mu(k),\end{align*}
are known. Two well-known examples, both of which may be found in Apostol (\cite[Theorem 2.1 and Theorem 3.12]{apostol1998introduction}), are
\begin{align}
\sum_{k|n}\mu(k)=\begin{cases} 1, &\text{if $n=1$;}\\
0, & \text{otherwise,}
\end{cases}
\label{eqn:delta}
\end{align}
and 
\begin{align}
\sum_{k\leq n}\floor{\frac{n}{k}} \mu(k)=1.
\label{eqn:sumdelta}
\end{align}
This pair of identities is closely related to Proposition~\ref{prp:qn} of Section~\ref{sec:qn}.
The Mertens function is also known to satisfy similar kinds of identities. One example is
\begin{align*}
\sum_{k\leq n} M\paren{\frac{n}{k}} &= 1. 
\end{align*}
Lehman \cite{lehman1960liouville}, Del{\'e}glise and Rivat \cite{deleglise1996computing} and Benito and Varona  \cite{benito2008recursive} have all leveraged 
this and similar identities to find efficient recursive methods of evaluating $M$ and related functions at large values.
We state a similar identity in Section~\ref{subsec:bv}.

\section{The functions  \texorpdfstring{$\epsilon^{m,n}$}{epsilon {m,n}} and \texorpdfstring{$E^{m,n}$}{E {m,n}}}
\subsection{Notation, statements of the main results, and proofs}
\label{sec:main}

The greatest common divisor of positive integers $a,b\in\N$ is written $(a,b)$.
Let $D$ be the matrix
\begin{align*}
D({i,j}) \coloneqq \begin{cases} 1, & \text{if $j|i$;}\\
0, & \text{otherwise.}
\end{cases}
\end{align*}
The action of this matrix on vectors corresponds to Dirichlet convolution by $\vv{1}\coloneqq (1,1,1,\ldots)$.
In matrix notation, where the functions $\mu$ and $\delta \coloneqq \paren{1,0,0,0,\ldots}$ are treated as a column vectors, 
and $*$ denotes Dirichlet convolution,
\begin{align}\vv{1}*\mu  = \paren{ \sum_{d|k} \mu(d) : k\geq 1 } = D\mu= \delta.\label{eqn:dirinv}\end{align}
We adopt the convention that the smallest index of sequences, vectors and matrices is $1$, and not $0$. Thus, 
the leading entry of $\delta$ is $\delta(1)=1$. 
Let $\delta_m$ be the column vector consisting of the 
$m$ leading terms of $\delta$ and let 
\begin{align}c^{m}\coloneqq (\delta_{2m}, -\delta_{2m}, \delta_{2m},-\delta_{2m},\ldots).\label{def:cm}\end{align}
Then $c^{m}$ is supported on integers $k\equiv \modd{1} {2m}$ and  $c^{m}$ is $4m$-periodic.
We indicate the transpose operation using a~$\phantom{}'$ symbol.
The functions of primary interest are now defined.

For each $m,n\geq 1$, set
\begin{align*}
\epsilon^{m,n}\coloneqq D'_nc^{m,n},
\end{align*}
where $D_n$ denotes the $n$th leading principal submatrix of $D$, $D'_n$ denotes the transpose of $D_n$ and $c^{m,n}$ consists of the 
$n$ leading terms of $c^{m}$.  
Written explicitly, for $1\leq j\leq n$,
\begin{align}
\epsilon^{m,n}(j) = \sum_{\substack{k\\jk\leq n}} c^{m}(jk).
\label{eqn:defcm}
\end{align}
For example, 
\begin{align*}
\epsilon^{1,6} &= 
D_{6}' c^{1,6}=
\begin{pmatrix}
 1& 1& 1& 1& 1& 1\\
 0& 1& 0& 1& 0& 1\\
 0& 0& 1& 0& 0& 1\\
 0& 0& 0& 1& 0& 0\\
 0& 0& 0& 0& 1& 0\\
 0& 0& 0& 0& 0& 1
\end{pmatrix}
\begin{pmatrix}
\phantom{-}  1 \\
\phantom{-}  0 \\
            -1 \\
\phantom{-}  0 \\
\phantom{-}  1 \\
\phantom{-}  0
\end{pmatrix}
= \begin{pmatrix}
\phantom{-}  1\\
\phantom{-}  0\\
            -1\\
\phantom{-}  0\\
\phantom{-}  1\\
\phantom{-}  0
\end{pmatrix}.
\end{align*}

The following proposition describes basic properties of  $\epsilon^{m,n}$.
When applied to the above example, items~\ref{item:bound} and \ref{item:identity} of Proposition~\ref{prp:main}
imply that $\norm{\epsilon^{1,6}}_{\infty}\leq 1$ and $\sum_{k\leq 6}\epsilon^{1,6}(k)\mu(k)=1$, which may be 
verified by inspection.
\begin{proposition}
In the established notation, 
\begin{enumerate}
\item \label{prp:gcd}if $1\leq k \leq n$ and $(k,2m)>1$, then $\epsilon^{m,n}(k)=0$;
\item \label{item:bound}$\epsilon^{m,n}\in\set{-1,0,1}^{n}$;
\item if $n\geq 3m$, then $\norm{\epsilon^{m,n}}_{1}\geq\floor{n/(3m)}$ and {$\max\!\set{k\!:\!\epsilon^{m,n}(k)\neq 0}\!=2m\!\floor{(n-1)/(2m)}\!+\!1$};\label{item:vanish}
\item \label{item:identity}$\sum_{k\leq n} \epsilon^{m,n}(k)\mu(k)=1$.
\end{enumerate}
\label{prp:main}
\end{proposition}
\begin{proof}
For all $1\leq j\leq n$,
\begin{align*}
\epsilon^{m,n}(j)&= 
\sum_{\substack{k\\jk\leq n}} c^{m}(jk)\\
&= \sum_{\substack{k\leq n\\k\equiv \modd{0} {j}\\k\equiv \modd{1} {4m}}} 1 - 
    \sum_{\substack{k\leq n\\k\equiv \modd{0} {j}\\ k\equiv \modd{2m+1} {4m}}} 1\\
&= \abs{A^{j}} - \abs{B^{j}},
\end{align*}
where \begin{align*}
A^{j}\coloneqq \set{  jk\leq n:  k\in\N,\, jk\equiv \modd{1} {4m}}\end{align*}
and 
\begin{align*}
B^{j}\coloneqq\set{ jk\leq n: k\in\N,\, jk\equiv \modd{2m+1} {4m}}. \end{align*}
To see item~\ref{prp:gcd} of the proposition, observe that neither $jk\equiv \modd{1} {4m}$ nor $jk\equiv \modd{2m+1} {4m}$ has an integer solution when $(j,2m)>1$.
Consequently, $\epsilon^{m,n}(j) = \abs{A^{j}} - \abs{B^{j}}=0-0=0$  when $(j,2m)>1$.

For item~\ref{item:bound}, observe that either $\epsilon^{m,n}(j)=0$, or else $(j,2m)=1$ and
\begin{align*}
\epsilon^{m,n}(j)=\abs{A^{j}}-\abs{B^{j}} =
\sum_{\substack{k\leq n/j\\ k\equiv \modd{j^*} {4m}}} 1 - 
\sum_{\substack{k\leq n/j\\ k\equiv \modd{(2m+1)j^*} {4m}}} 1,
\end{align*}
where $j^*$ is the unique integer satisfying both $1\leq j^*\leq 4m-1$ and $jj^*\equiv \modd{1} {4m}$. Then since $j^*$ is necessarily
odd, $(2m+1)j^*\equiv \modd{2m+j^*} {4m}$. Therefore,
\begin{align*}
\epsilon^{m,n}(j)&= 
\sum_{\substack{k\\ 0<j^* + 4mk\leq n/j}}1 - \sum_{\substack{k\\ 0<2m + j^* + 4mk\leq n/j}}1 \\
&=\sum_{0<2k+\frac{j^*}{2m}\leq \frac{n/j}{2m}}1 - \sum_{ 0<(2k+1) + \frac{j^*}{4m}\leq \frac{n/j}{2m}} 1.
\end{align*}
Hence,
\begin{align}
\epsilon^{m,n}(j)&= \sum_{-\frac{j^*}{2m}<v\leq \frac{n/j}{2m}-\frac{j^*}{2m}} \paren{-1}^v.
\label{eqn:explicit}
\end{align}
Consecutive terms in the last expression sum to 0. Since each individual term in the expression is either $-1$ or $+1$, 
the total sum is either $-1$, $0$, or $1$.

We now show the third claim. 
Assume $n\geq k>n/3$.
From Eq.~(\ref{eqn:defcm}), and
since $c^{m}(k)=0$ whenever $k\not\equiv \modd{1} {2m}$,
\begin{align}
\epsilon^{m,n}(k) &= c^{m}(k)+c^{m}(2k) = c^{m}(k).
\label{eqn:cm}
\end{align}
Consequently,
\begin{align*}
\norm{\epsilon^{m,n}}_1 \geq \sum_{n/3<j\leq n} \abs{\epsilon^{m,n}(j)} = \sum_{n/3<j\leq n} \abs{c^{m}(j)} &=  \sum_{\substack{n/3<j\leq n\\j \equiv \modd{1} {2m}}} 1
= \sum_{n/3<1+2ml\leq n} 1
= \sum_{\frac{n-3}{6m}< l \leq \frac{3n-3}{6m} } 1.
\end{align*}
Since $a<b$ implies $\abs{ (b-a) - \sum_{a<j\leq b}1}<1$, it follows that
\begin{align*}
\norm{\epsilon^{m,n}}_1  > \frac{n}{3m}-1.
\end{align*}
Since $\norm{\epsilon^{m,n}}_1 $ is integer-valued, the first statement of the third item follows. The second statement of this item follows from the definition of $c^{m}$ and Eq.~(\ref{eqn:cm}).

To see the last claim, observe that Eq.~(\ref{eqn:dirinv}) implies $D_n\mu_n=\delta_n$, where $\mu_n$ represents the leading $n$ terms of $\mu$.
Then, by applying the definitions of $c^{m,n}$ and $\epsilon^{m,n}$, it follows that
\begin{align*}
\sum_{1\leq k\leq n} \epsilon^{m,n}(k)\mu(k) &=  \mu_n'\epsilon^{m,n} =  \mu_n'D_n' c^{m,n} =  \paren{D_{n}\mu_n}'  c^{m,n} =  \delta_n' c^{m,n} =  c^{m,n}(1) = 1.\qedhere
\end{align*}
\end{proof}
Equation~(\ref{eqn:cos}) of Section~\ref{sec:intro} follows from the above observations. 
To see this, let $d^{k,n}$ denote column $k$ of $D_n$ and note that $d^{k,n}$ is supported on integers $jk\leq n$.
Next, observe that
\begin{align*}
c^{1,n} &= (1,0,-1,0,\ldots) = \paren{\cos{\paren{\frac{\pi (k-1)}{2}}}:1\leq k\leq n}.
\end{align*}
Hence,
\begin{align*}
\epsilon^{1,n}(k) = \paren{d^{k,n}}' c^{1,n} = \sum_{jk\leq n}\cos{\paren{\frac{\pi (jk - 1)}{2}}},
\end{align*}
which shows the claim. 

\begin{corollary}
Let $m,n,\in\N$. For each positive integer $k\leq n$ set \begin{align*}
\omega^{m,n}(k)\coloneqq \begin{cases}
\epsilon^{m,n}(k), & \text{if $k$  is odd;}\\
\epsilon^{m,{\floor{n/2}}}(k/2), & \text{if $k$ is even.}
\end{cases}
\end{align*}
Then $\omega^{m,n}\in \set{-1,0,1}^{n}$. If $n>1$, then $\sum_{k\leq n}\omega^{m,n}(k)\mu(k)=0$.
\label{cor:omega}
\end{corollary}
\begin{proof}
By Proposition~\ref{prp:main} item \ref{prp:gcd},
for all $m,n\geq 1$, the function $\epsilon^{m,n}$ is supported on odd integers.  
Since $\mu(2k)=-\mu(k)$ when $k$ is odd,  it follows by Proposition~\ref{prp:main} item \ref{item:identity} that
\begin{align*}
\sum_{k\leq n}\omega^{m,n}(k)\mu(k) &= 
\sum_{2k+1\leq n}\epsilon^{m,n}(2k+1)\mu(2k+1) + 
\sum_{2k\leq n}\epsilon^{m,{\floor{n/2}}}(k)\mu(2k) = 1-1=0.\qedhere
\end{align*}
\end{proof}

Our next aim is to develop analogous results for 
the summatory function of $\mu$, {i.e.}, the Mertens function, $M$.
Let $S_{n}\in\set{0,1}^{n\times n}$ denote the lower-triangular matrix with $S_n(i,j)=1$ whenever $j\leq i$.
Then $S_n$ is the matrix representation
that corresponds to the operation of cumulative summation applied to $n$-vectors.
Let $M_n\coloneqq (M(1),\ldots,M(n))$. Then $M_n = S_n\mu_n$.
Let $\Delta_n\coloneqq \paren{S_{n}'}^{-1}$. Then 
\begin{align}
\Delta_n &= \begin{pmatrix}
1 &          -1 \\
  &\phantom{-}1 & -1\\
  &             & \phantom{-}1 & -1\\
  &             &  & \ddots &-1\\
  &             &  &   &\phantom{-}1
\end{pmatrix},
\label{eqn:deltan}
\end{align}
which is the matrix representation of a difference operation applied to $n$-vectors.  Set 
\begin{align*}
E^{m,n}\coloneqq \Delta_n^{\phantom{m}}\epsilon^{m,n}. 
\end{align*}
The next corollary combines the above observations into a result that parallels Proposition~\ref{prp:main}.
\begin{corollary}
In the established notation,
\begin{enumerate}
\item $E^{m,n}\in\set{-1,0,1}^n$;
\item if $n\geq 3m$, then $\norm{E^{m,n}}_1\geq 2\floor{n/(3m)}-1$;
\item $\sum_{k\leq n} M(k) E^{m,n}(k)=1$.
\end{enumerate}
\label{cor:main}
\end{corollary}
\begin{proof}
Since $\epsilon^{m,n}(2k)=0$ for all $k\geq 1$, the difference of adjacent entries of 
$\epsilon^{m,n}$ is $-1,0$ or~$1$. This shows the first claim.

For $1\leq k<n$,
\begin{align*}
\abs{E^{m,n}(k)} = \abs{\epsilon^{m,n}(k)} + \abs{\epsilon^{m,n}(k+1)}.
\end{align*}
This fact, when combined with items~\ref{item:bound} and \ref{item:vanish} of Proposition~\ref{prp:main} implies
\begin{align*}
\norm{E^{m,n}}_1 &= \sum_{k<n}\paren{ \abs{\epsilon^{m,n}(k)} + \abs{\epsilon^{m,n}(k+1)}} + \abs{\epsilon^{m,n}(n)}\\
&=\norm{\epsilon^{m,n}}_1 + \norm{\epsilon^{m,n}}_1 -\abs{\epsilon^{m,n}(1)} 
\geq 2\floor{\frac{n}{3m}}-1.
\end{align*}
This shows the second claim.

Finally, apply Proposition~\ref{prp:main} and the respective definitions to find
\begin{align*}
M'_n E^{m,n}  &= \paren{S_n\mu_n}' \Delta^{\phantom{m}}_n\epsilon^{m,n}
= \mu_n' S_n' \paren{S_{n}'}^{-1}\epsilon^{m,n}
= \mu_n'\epsilon^{m,n}
= 1.
\qedhere\end{align*}
\end{proof}

\subsection{Remarks  on Eq.~(\ref{eqn:cos})}
\label{subsec:cos}
It follows from (\ref{eqn:cos}), that for all positive integers $n$ and $k\in\N$ satisfying $k\leq n$, 
the value of $\epsilon^{1,n}(k)$ equals the real part of the sum, ${\sum_{1\leq j\leq \floor{n/k}} i^{jk-1}}$,
where $i\coloneqq \sqrt{-1}$.
This leads to a closed-form expression for $\epsilon^{1,n}(k)$. When $k$ is even, 
it follows by Proposition~\ref{prp:main} item \ref{prp:gcd} that $\epsilon^{1,n}(k)=0$. Now assume $k$ is odd. One has
\begin{align*}
\sum_{1\leq j\leq \floor{n/k}} i^{jk-1} & = i^{k-1}\frac{1 - i^{k\floor{n/k}}}{1-i^{k}} = 
(-1)^{\floor{k/2}}\frac{(1 - i^{k\floor{n/k}})(1+i^{k})}{2}.
\end{align*}
Since $i^{k\floor{n/k}}=\paren{i^{\floor{n/k}}}^{k}$, the above expression vanishes 
when $\floor{n/k}\equiv \modd{0} {4}$.
Next, when $\floor{n/k}$ is odd, and since $k$ is assumed odd, the real part of the above expression
equals $(-1)^{\floor{k/2}}\paren{ 1 - i^{k(\floor{n/k}+1)}}/2$, which vanishes when $\floor{n/k}\equiv \modd{-1} {4}$.
When $\floor{n/k}\equiv \modd{1} {4}$, one has $i^{\floor{n/k}}=i$ so that the right-hand side of the above 
sum reduces to 
\begin{align*}
(-1)^{\floor{k/2}} \frac{(1-i^{k})(1+i^{k})}{2} &= (-1)^{\floor{k/2}} \frac{1-i^{2k}}{2} = (-1)^{\floor{k/2}}.
\end{align*}
Finally, when  $\floor{n/k}\equiv \modd{2} {4}$, $i^{\floor{n/k}}=-1$, and (recalling that $k$ is assumed odd) the right-hand side of the above sum reduces to
\begin{align*}
(-1)^{\floor{k/2}} \frac{(1-(-1)^{k})(1+i^{k})}{2} &= (-1)^{\floor{k/2}}(1+i^{k}),
\end{align*}
the real part of which is $(-1)^{\floor{k/2}}$.

Therefore, for all positive integers $n$ and $k$ that satisfy $k\leq n$,
\begin{align}
\epsilon^{1,n}(k)&= \begin{cases} 0, & \text{if $\floor{n/k}\equiv \modd{r} {4}$ for some $r\in\set{0,-1}$};\\
      \sin\paren{k\pi/2},&\text{otherwise.}
\end{cases}
\label{eqn:e1n}
\end{align}

\label{subsec:explicit}
When $1\leq k\leq n$, an explicit formula for $\epsilon^{m,n}(k)$ is possible to derive. 
The case where $(k,2m)>1$ is discussed in Proposition~\ref{prp:main} item~\ref{prp:gcd}, so we now assume $(k,2m)=1$.
Let $k^*$ be any integer satisfying $kk^*=\modd{1} {4m}$.  By Eq.~(\ref{eqn:explicit}), one has
\begin{align*}
\epsilon^{m,n}(k)&=
\sum_{\substack{j\\
\frac{-k^*}{2m}<2j\leq \frac{\floor{n/k} - k^*}{2m}}} 1
 -
\sum_{\substack{j\\
\frac{-k^*}{2m}<2j+1\leq \frac{\floor{n/k} - k^*}{2m}}} 1\\
&=
\sum_{\substack{j\\
\frac{-k^*}{4m}<j\leq \frac{\floor{n/k} - k^*}{4m}}} 1
 -
\sum_{\substack{j\\
\frac{-k^*}{4m}-\frac{1}{2}<j\leq \frac{\floor{n/k} - k^*}{4m}-\frac{1}{2}}} 1\\
&=
\paren{\floor{\frac{\floor{n/k} - k^*}{4m}} - \floor{\frac{-k^*}{4m}}}
-
\paren{\floor{\frac{\floor{n/k} - k^*}{4m}-\frac{1}{2}} - \floor{\frac{-k^*}{4m}-\frac{1}{2}}}\\
&=
\paren{\floor{\frac{\floor{n/k} - k^*}{4m}} - \floor{\frac{\floor{n/k} - k^*}{4m}-\frac{1}{2}} }
-
\paren{
\floor{\frac{-k^*}{4m}}
- \floor{\frac{-k^*}{4m}-\frac{1}{2}}}\\
&= 
\sum_{\substack{j\\
\frac{\floor{n/k}-k^*}{4m}-\frac{1}{2} < j\leq {\frac{\floor{n/k} - k^*}{4m}}}} 1
- 
\sum_{\substack{j\\
\frac{- k^*}{4m}-\frac{1}{2}< j \leq {\frac{- k^*}{4m}}}} 1.
\end{align*}
Then, letting $\pbrace{x}\coloneqq x-\floor{x}$ denote the fractional part of $x$,
\begin{align*}
\epsilon^{m,n}(k)&=\begin{cases}
\phantom{-} 1, &\text{if $(k,2m)=1$, $\pbrace{\frac{k^*}{4m}}\in (0,\frac{1}{2}]$ and $\pbrace{\frac{k^*-\floor{n/k}}{4m}}\not\in (0,\frac{1}{2}]$;}\\
           -1, &\text{if $(k,2m)=1$, $\pbrace{\frac{k^*}{4m}}\not\in (0,\frac{1}{2}]$ and $\pbrace{\frac{k^*-\floor{n/k}}{4m}}\in (0,\frac{1}{2}]$;}\\
\phantom{-} 0, &\text{otherwise.}
\end{cases}
\end{align*}
\subsection{An identity involving a sparse sum}
\label{subsec:bv}
Benito and Varona (\cite[Theorem 9]{benito2008recursive}) establish the identity,
\begin{align*}
2M(n)+3 = g(n,1) + \sum_{3\leq a\leq n-1} h(a)\paren{ M\paren{\frac{n}{a}} - M\paren{\frac{n}{a+1}}},
\end{align*}
where $g(x,1)$ and $h(x)$ depend on the value of $\modd{x} {6}$, and the 
values of both $g$ and $h$ are tabulated. 
The utility of this formula is that it leads to an efficient and recursive means of evaluating $M$.
In Proposition~\ref{cor:bv} of this section, we establish a similar identity. 

Before stating the identity, we introduce notation.
The unique non-principal Dirichlet character mod 4 is
\begin{align*}
\chi(k)&=\begin{cases} 
\phantom{-}1, &\text{if $k\equiv \modd{1} {4}$;}\\
-1, &\text{if $k\equiv \modd{-1} {4}$;}\\
\phantom{-}0, &\text{otherwise}.
\end{cases}
\end{align*}
The sine-factor in the non-vanishing case of~(\ref{eqn:e1n}) satisfies 
$\sin\paren{k\pi/2} = \chi(k)$ for all integers $k$.
Next, we define 
\begin{align*}
M(x,\chi)\coloneqq \sum_{k\leq x} \mu(k)\chi(k),
\end{align*}
and set
\begin{align*}
\eta(k)\coloneqq \begin{cases} 1, &\text{if $k\equiv \modd{r} {4}$ for some $r\in\set{1,2}$;}\\
0, &\text{otherwise.}
\end{cases}
\end{align*}
We are now ready to state the identity. 
\begin{proposition}
For all $n\geq 1$,
\begin{align*}
\sum_{1\leq k\leq n} \eta(k)\paren{ M\paren{\frac{n}{k},\chi} -  M\paren{\frac{n}{k+1},\chi} } = 1.
\end{align*}
\label{cor:bv}
\end{proposition}
\begin{proof}
One has
\begin{align*}
\sum_{1\leq k\leq n} \eta(k)\paren{ M\paren{\frac{n}{k},\chi} -  M\paren{\frac{n}{k+1},\chi} } &=
\sum_{\substack{k,j\\k\equiv \modd{r} {4}, r\in\set{1,2}\\{\frac{n}{k+1}< j\leq \frac{n}{k}}}} \mu(j)\chi(j)\\
&=\sum_{4l+3>\floor{n/j}\geq 4l+1} \mu(j)\chi(j)\\
&=\sum_{j} \epsilon^{1,n}(j)\mu(j)\\
&=1,
\end{align*}
where the third equality follows from Eq.~(\ref{eqn:e1n}) and the definition of $\chi$, 
and the last equality from Proposition~\ref{prp:main}~item~\ref{item:identity}.
\end{proof}

\section{A family of unimodular matrices}
\label{sec:qn}
The classical identities expressed in Eq.~(\ref{eqn:delta}) and Eq.~(\ref{eqn:sumdelta}) of the introduction are equivalent.
Indeed, for any integer $n\in\N$, 
Eq.~(\ref{eqn:delta}) states $D_n\mu_n=\delta_n$. 
Multiplication of both sides by $S_n$ yields, $S_nD_n\mu_n=S_n\delta_n = \vv{1}_n$, which is (\ref{eqn:sumdelta}). 
Equivalence follows by observing that $S_n$ is invertible. This observation generalizes to 
yield the infinite family of identities, $S_n^{k}D^{\phantom{k}}_n\mu_n = S_n^{k}\delta^{\phantom{k}}_n$, where
$k,n\in\N$ are arbitrary.  The next proposition describes a closely-related family of identities.
We first state a lemma that generalizes the identity, $\sum_{d|n}\phi(d)=n$, where $\phi$ is Euler's totient function.
\begin{lemma}
For positive $j,k\in\N$, write $\phi^{k}(j)\coloneqq \sum_{\ell\leq j/k,\, (\ell,j)=1} 1$. 
Let $R_{n}\coloneqq D_n^{-1}S_nD_n$. Then for all $1\leq j,k\leq n$, $R_{n}(j,k) = \phi^{k}(j)$. 
\label{lem:phi}
\end{lemma}
\begin{proof}
Column $k$ of $S_nD_n$ equals $(\floor{1/k},\floor{2/k},\ldots,\floor{n/k})$. Let $\phi^{k}_{n}$ denote the column vector consisting of the 
$n$ leading terms of $\phi^{k}$.  Our aim is to demonstrate that
\begin{align}
(D_n\phi^{k}_{n})(j) &= \sum_{d|j}\phi^{k}(d)=\floor{j/k}
\label{eqn:phin}
\end{align}
holds when $1\leq j,k \leq n$. Once this has been shown, one arrives at $( D_n^{-1}S_nD_n)(j,k)=\phi^{k}(j)$ by applying $D_n^{-1}$ to both sides.
Fix $j$ and $k$. For each integer $d$ with $d|j$, define the set
\begin{align*}
S_d\coloneqq\set{s\in\N: 1\leq s\leq j/k\text{ and } (s,j)=j/d}.
\end{align*}
Then $S_d$ consists of all integers with the form $(t)(j/d)$, where $1\leq t\leq d/k$ and $(t,d)=1$.
Then $\abs{S_d} = \phi^k(d)$. Each integer between $1$ and $\floor{j/k}$ belongs to a unique $S_d$. Summing over $d$, one has~(\ref{eqn:phin}).
\end{proof}
\begin{proposition}
For positive $n\in\N$, set $Q_n\coloneqq R_n'R_n$.
Then $Q_n$ is a positive definite matrix, both $Q_n$ and $(Q_n)^{-1}\in\Z^{n\times n}$, and $\det Q_{n} = \det\paren{Q_{n}}^{-1}=1$. 
Additionally,
\begin{enumerate}
\item $Q_n\mu_n=\delta_n$\label{item:delta}; 
\item $Q_{n}(i,j)>0$ for all $1\leq i,j\leq n$;\label{item:pos}
\item $Q_{n}(1,1) = \sum_{1\leq k\leq n} \phi(k)^2$\label{item:lambdan}. 
\end{enumerate}
\label{prp:qn}
\end{proposition}
\begin{proof}
Recall that $D$ is the matrix representation of Dirichlet convolution by $\vv{1}$, and the inverse operation is Dirichlet convolution by $\mu$.
In particular, the matrix $D_n^{-1}$ is $(-1,0,1)$-valued.
Since each of $D_n$, $D_n^{-1}$ and $S_n$ live in $\Z^{n\times n}$, it follows that $Q_n\in\Z^{n\times n}$. 

Also, each of $D_n$, $D_n^{-1}$ and $S_n$ are lower-triangular with $1$ along the main diagonal, so each has determinant equal to $1$, as
do their transposes. Since
$\det$ is multiplicative, $\det Q_n=1$, from which it follows (by Cramer's rule) that $(Q_n)^{-1}\in\Z^{n\times n}$. Additionally, since for all $x\in\R^{n}\setminus\set{\vv{0}}$ one has $x'Q_nx = \norm{(D_n^{-1}S_nD_n) x}^2\geq 0$,
$Q_n$ is nonnegative definite. But $\det Q_n=1>0$, so it has full rank. In particular, $x'Q_nx >0$, which implies, since $x\neq \vv{0}$ is arbitrary, that $Q_n$ is positive definite.

We now establish item~\ref{item:delta} of the proposition's statement.
The first column of $S_n$ equals the first column of $D_n$, both of which equal $\vv{1}_n$. Thus, recalling Eq.~(\ref{eqn:dirinv}),
\begin{align*}
R_{n}\mu_n = D_n^{-1}S_nD_n\mu_n = D_n^{-1}S_n\delta_n = D_n^{-1}\vv{1}_n=\delta_n.
\end{align*}
The first column of each of the upper-triangular matrices $(D_n^{-1})'$, $D_n'$ and $S_n'$ equals $\delta_n$, so
\begin{align*}
Q_n\mu_n = R_n'R_n\mu_n = D_n'S_n'\paren{D_n^{-1}}'\delta_n = D_n'S_n'\delta_n = D_n'\delta_n  =\delta_n.
\end{align*}

For the proof of item~\ref{item:pos}, note that $\phi^k$ of Lemma~\ref{lem:phi} satisfies $\phi^{k}(l)>0$ whenever $l\geq k$. 
Therefore $Q_{n}(j,k)=\sum_{l\leq n} \phi^{j}(l)\phi^k(l)\geq \phi^{j}(n)\phi^{k}(n)>0$ when $1\leq j,k\leq n$.
Item~\ref{item:lambdan} is also immediate from Lemma~\ref{lem:phi}. Indeed, $Q_{n}(1,1)=\sum_{j\leq n} \paren{\phi^1(j)}^{2}=\sum_{j\leq n}(\phi(j))^{2}$.
\end{proof}


In the discussion below, it may be helpful to have a reference example of $R_{n}$, $Q_n$ and~$Q_{n}^{-1}$:
\begin{align*}
R_7 = \begin{pmatrix}
1 \\
1 & 1\\
2 & 1 & 1 \\
2 & 1 & 1 & 1\\
4 & 2 & 1 & 1 & 1\\
2 & 1 & 1 & 1 & 1 & 1\\
6 & 3 & 2 & 1 & 1 & 1 & 1
\end{pmatrix},\qquad
Q_7 = \begin{pmatrix}
  66 & 33 &  22  & 14  & 12   & 8   & 6\\ 
  33 & 17 &  11  &  7  &  6   & 4   & 3\\ 
  22 & 11 &   8  &  5  &  4   & 3   & 2\\ 
  14 &  7 &   5  &  4  &  3   & 2   & 1\\ 
  12 &  6 &   4  &  3  &  3   & 2   & 1\\ 
   8 &  4 &   3  &  2  &  2   & 2   & 1\\ 
   6 &  3 &   2  &  1  &  1   & 1   & 1  
\end{pmatrix},
\end{align*}
and
\begin{align*}
(Q_7)^{-1} = \begin{pmatrix}
\hm   1 &     -1  &     -1 & \hm    &     -1  & \hm  1  &      -1\\ 
     -1 & \hm  2  & \hm    & \hm    & \hm     & \hm     & \hm    \\ 
     -1 & \hm     & \hm  3 &     -1 & \hm  2  &     -2  & \hm   1\\ 
\hm     & \hm     &     -1 & \hm  2 &     -1  & \hm     & \hm   1\\ 
     -1 & \hm     & \hm  2 &     -1 & \hm  4  &     -3  & \hm   2\\ 
\hm   1 & \hm     &     -2 & \hm    &     -3  & \hm  4  &      -3\\ 
     -1 & \hm     & \hm  1 & \hm  1 & \hm  2  &     -3  & \hm   5  
\end{pmatrix}.
\end{align*}

Item \ref{item:delta} of the proposition states that $\mu_n$ is the first column of $(Q_{n})^{-1}$. The second column of $(Q_{n})^{-1}$ has a far simpler form, namely $q\coloneqq(-1,2,0,0,0,\ldots,0)$.
To see this, observe that $R_n(2,1)=R_n(2,2)=1$. When $j>2$, $R_n(j,2)=\floor{R_n(j,1)/2}$, and
since $R_n(j,1)=\phi^{1}(j)=\phi(j)$ is even, $\floor{R_n(j,1)/2}={R_n(j,1)/2}$.
Noting that $R_n'$ is upper-triangular and $R(1,1)=1$, one has 
\begin{align*}R_n'R_nq=R_n'\begin{pmatrix}-1\\\hm 1\\\hm 0\\\hm\vdots\\\hm 0\end{pmatrix} 
= \begin{pmatrix}0\\ 1\\ 0\\\vdots\\0\end{pmatrix}.\end{align*}

Since $Q_n$ is Hermitian, the eigenvalues of $Q_n$  are real.  Let 
$\lambda_1\leq \lambda_2\leq \ldots\leq \lambda_n$ denote the eigenvalues of $Q_n$.
By item~\ref{item:pos}, we may invoke the Perron-Frobenius theorem to conclude that $\lambda_n$ is simple and that the eigenvector associated
with $\lambda_n$, denoted $v_n$, satisfies $v_n(j)>0$ for $1\leq j\leq n$. 

The following two propositions concern the spectrum of $Q_n$. Table~\ref{tbl:lambdai} illustrates these propositions.
\begin{table}
\begin{center}
\caption{Eigenvalue $\lambda_i$ of $Q_n$.}
\label{tbl:lambdai}
\begin{tabular}{clll}
\toprule
$i$ & $n=999$ & $n=1000$ & $n=1001$\\
\midrule
1           &  $ 4.588 \times 10^{-4}$ &   $4.580   \times 10^{-4}$ &  $   4.573 \times 10^{-4}$\\
2           &  $ 6.887 \times 10^{-4}$ &   $6.879   \times 10^{-4}$ &  $   6.859 \times 10^{-4}$\\
$\cdots$\\
$n-1$       &  $ 4.458 \times 10 ^{4}$ &   $4.471\times 10^{4}$ &     $4.487\times 10^{4}$\\
$n$         &  $ 2.344 \times 10 ^{8}$ &   $2.346\times 10^{8}$ &     $2.355\times 10^{8}$\\
\bottomrule
\end{tabular}
\end{center}
\end{table}

\begin{proposition}
The eigenvalues of $Q_n$ interlace the eigenvalues of $Q_{n+1}$. 
That is, if $\lambda_{i}$ is the $i$th eigenvalue of $Q_n$ and $\gamma_{i}$ the $i$th eigenvalue of $Q_{n+1}$, then
\begin{align*} \gamma_{1}\leq \lambda_{1}\leq \gamma_{2}\leq \cdots \leq \gamma_{n}\leq \lambda_n\leq \gamma_{n+1}.\end{align*}
\end{proposition}
\begin{proof}
We claim that $(Q_{n})^{-1}$ is a leading principal submatrix of $(Q_{n+1})^{-1}$. 
Assume, for the moment, that this has been established.  Then by Cauchy's interlacing theorem, 
the eigenvalues of $(Q_{n})^{-1}$ interlace those of $(Q_{n+1})^{-1}$. That is, letting $\lambda_{i}^{-1}$ and $\gamma_{i}^{-1}$ denote the eigenvalues 
of $(Q_{n})^{-1}$ and $(Q_{n+1})^{-1}$, respectively, 
\begin{align*} \gamma_{n+1}^{-1}\leq \lambda_{n}^{-1}\leq \gamma_{n}^{-1} \leq \cdots \leq  \gamma_{2}^{-1}\leq \lambda_{1}^{-1}\leq \gamma_{1}^{-1}.\end{align*}
Since all quantities in this expression are positive scalars, the proposition follows.

We now show that $(Q_{n})^{-1}$ is a leading principal submatrix of $(Q_{n+1})^{-1}$.  
From the definition of~$Q_n$, we have that
\begin{align}
Q_{n}^{-1} = (R_{n}'R_{n})^{-1} =(R_{n})^{-1}(R_{n}')^{-1} = (R_{n})^{-1}(R_{n}^{-1})'.\label{eqn:qnrn}
\end{align} 
Recall from Lemma~\ref{lem:phi} the entrywise equality, $R_{n}(j,k)=\phi^{k}(j)$. Since the right-hand side does not depend on $n$, the $(j,k)$ entry 
of $R_{n}$ does not depend on $n$. Therefore, $R_{n}$ is a leading principal submatrix of $R_{m}$ whenever $m\geq n$.
Since $R_{n}$ is lower-triangular, it follows (by consideration of block matrix inversion, for example) that $(R_{n})^{-1}$ is a 
leading principal submatrix of $(R_{m})^{-1}$ for all $m\geq n$.
This also shows that $(R_{n}^{-1})'$ is a leading principal submatrix of $(R_{m}^{-1})'$ whenever $m\geq n$. 

For $1\leq j\leq n$, define $r^j$ as the $j$th column of $(R_{n+1}^{-1})'$. 
Since $(R_{n}^{-1})'$ is a leading principal submatrix of $(R_{n+1}^{-1})'$, the first $n$ entries of $r^j$ agree with the $j$th column of $(R_{n}^{-1})'$.
This, combined with Eq.~(\ref{eqn:qnrn}) and the fact that $R_{n}^{-1}$ is lower-triangular, implies that the 
$j$th column of $(Q_{n})^{-1}$ equals the column vector that is comprised 
of the first $n$ entries of $(R_{n+1})^{-1} r^j$. 
Eq.~(\ref{eqn:qnrn}) also implies that the $j$th column of $(Q_{n+1})^{-1}$ equals $(R_{n+1})^{-1} r^j$. This applies to all 
$1\leq j\leq n$, so it follows that $(Q_{n})^{-1}$ is a leading principal submatrix of  $(Q_{n+1})^{-1}$.
\end{proof}

The next proposition establishes bounds on the extremal eigenvalues of $Q_n$.
\begin{proposition}
Let $\lambda_{n}$ denote the dominant eigenvalue of $Q_n$. Then
\begin{align*}
\paren{c  + o(1)}n^{3}\leq \lambda_n \leq {\frac{\pi^2}{18}n^{3} + O(n^{2})},
\end{align*}
where $c$ is a constant that is approximately equal to $0.1427$.
Let $\lambda_1$ denote the smallest eigenvalue of $Q_n$. Then
\begin{align*}
\lambda_1\leq {\frac{\pi^2}{6}n^{-1} + O\paren{n^{-3/2}}}.
\end{align*}
\label{prp:lambdanbound}
\end{proposition}
\begin{proof}
We first show the lower bound on $\lambda_n$.
Since $\lambda_n$ is the dominant eigenvalue of $Q_n$, applying item~\ref{item:lambdan} of Proposition~\ref{prp:qn} one has,
\begin{align*}
\lambda_n \geq \frac{\delta_n' Q_n \delta_n}{\norm{\delta_n}^2} = Q_n(1,1) = \sum_{k\leq n} \paren{\phi(k)}^2.
\end{align*}
We estimate the right-hand sum by applying the Wiener-Ikehara Theorem to the function $H(s)\coloneqq\sum_{k\geq 1}\phi(k)^2 k^{-s}$.
A comment associated with the sequence~\seqnum{A127473} asserts that $H(s)=\zeta(s-2)\prod_{p\ \textrm{prime}}\paren{1-2p^{1-s} + p^{-s}}$.
Since this is a known fact, but no proof is provided in the citation, we sketch a proof here. To start, if $n=\prod_{p} p^{r}$, then
$\phi(n)=\prod_{p|n} p^{r}-p^{r-1}$. Since $n\mapsto \phi(n)n^{-s}$ is a multiplicative function, one has the Euler product,
\begin{align*}
H(s)&=\prod_{p\ \textrm{prime}} \paren{1 + \sum_{r\geq 1} \frac{\paren{p^{r}-p^{r-1}}^2}{p^{rs}}}\\
&= \prod_{p\ \textrm{prime}} \paren{1 + \sum_{r\geq 1} \frac{p^{2r}}{p^{rs}} - 2\frac{p^{2r-1}}{p^{rs}} + \frac{p^{2r-2}}{p^{rs}}}\\
&= \prod_{p\ \textrm{prime}} \paren{1 + \paren{1 - \frac{2}{p} + \frac{1}{p^{2}}}\frac{p^{2-s}}{1-p^{2-s}}}\\
&= \prod_{p\ \textrm{prime}} \paren{1-p^{2-s} + {p^{2-s} - {2}{p^{1-s}} + {p^{-s}}}} \paren{\frac{1}{1-p^{2-s}}}\\
&= \zeta(s-2) \prod_{p\ \textrm{prime}} \paren{1 - 2p^{1-s}+p^{-s}}.
\end{align*}
At $s=3$, the infinite product converges and $\zeta(s-2)$ has a simple pole with residue equal to 1. Therefore,
$H$ has a simple pole at $s=3$, with residue equal to the product of the right-hand side.
We now apply the Wiener-Ikehara Theorem to conclude that
\begin{align*}
\sum_{k\leq n}\paren{\phi(x)}^2 \sim \frac{n^3}{3} \prod_{p\ \textrm{prime}}  \paren{1 - 2p^{-2}+p^{-3}}=\paren{\frac{n^3}{3}}\paren{ 0.42824950\ldots}
\end{align*}
We now show the upper bound on $\lambda_n$. Indeed,
\begin{align*}
\operatorname{Trace}\paren{Q_n} &= \sum_{k=1}^{n}Q_n(k,k) \\
&= \sum_{k=1}^{n}\sum_{j=1}^{n} \paren{\phi^{k}(j)}^2 \\
&\leq \sum_{k=1}^{\infty}\sum_{j=1}^{n}\paren{\frac{j}{k}}^2 \\
&= \frac{n(n+1)(2n+1)}{6} \sum_{k=1}^{\infty} k^{-2} \\
&= \frac{\pi^2}{18}n^3 + O(n^{2}),
\end{align*}
where we have used the elementary identity $\sum_{j=1}^n j^2=n(n+1)(2n+1)/{6}$ and Euler's identity $\sum_{i=1}^{\infty}{i^{-2}}=\pi^2/6$.
Finally, recall that $Q_n$ is positive definite and $\operatorname{Trace}\paren{Q_n} =\sum_{i=1}^{n}\lambda_i\geq \lambda_n$.

We now show the upper bound on $\lambda_1$. We have by item \ref{item:delta} of Proposition~\ref{prp:qn} that
\begin{align*}
\frac{\mu_n'Q_n\mu_n}{\norm{\mu_n}^2} = \frac{\mu_n'\delta_n}{\norm{\mu_n}^2} = \frac{\mu(1)}{\sum_{k\leq n}\mu_n(k)^2}
=
\frac{1}{\sum_{k\leq n}\abs{\mu_n(k)}}\geq \lambda_1.
\end{align*}
The penultimate expression is the reciprocal of the number of squarefree integers not exceeding~$n$. The desired bound is a consequence of this,
combined with an application of Gegenbauer's estimate \cite[p.\ 47]{gegenbauer1885},
\begin{align*}
{\sum_{k\leq n}\abs{\mu_n(k)}}=\frac{6}{\pi^2} n + O(n^{1/2}). 
\end{align*}
That is,
\begin{align*}
\lambda_1 &\leq \paren{\sum_{k\leq n}\abs{\mu_n(k)}}^{-1}={\frac{\pi^2}{6}n^{-1} + O(n^{-3/2})}.\qedhere
\end{align*}
\end{proof}
\begin{corollary}
In the established notation,
$\abs{v_n'\mu_n} \leq n^{-3/2} C(1+o(1))$,
\label{cor:mun}
where the constant $C$ is approximately 2.6467.
\end{corollary}
\begin{proof}
For $1\leq i\leq n$, let $v_i$ satisfy $Q_nv_i = \lambda_iv_i$. We may assume that $v_i'v_j=0$ whenever $i\not=j$, that $v_i'v_i=1$ for all $1\leq i\leq n$ and that the space spanned
by $\set{v_i:1\leq i\leq n}$ has dimension $n$.  Then $\mu_n = \sum_{1\leq i\leq n}(v_i'\mu_n) v_i$. By item~\ref{item:delta} of Proposition~\ref{prp:qn}, $\mu_n'Q_n\mu_n=1$. Since
$\lambda_i>0$ for all $i$,
\begin{align*}
1  = \mu_n'Q_n\mu_n = \sum_{1\leq i\leq n} \lambda_i \abs{ v_{i}'\mu_n}^2 \geq \lambda_n \abs{v_n'\mu_n}^2 \geq n^{3}(c+o(1)) \abs{v_n'\mu_n}^2,
\end{align*}
where the last inequality is from Proposition~\ref{prp:lambdanbound}.  The constant $C$ in the Corollary statement is $c^{-1/2}$.
\end{proof}

We conclude with a conjecture.
Recall that the dominant eigenvector of $Q_n$, which is denoted $v_n$, was shown to satisfy $v_n(j)>0$, for all $1\leq j\leq n$.
Assume we have normalized~$v_n$, so that $\norm{v_n}_2=1$. Let $h(k)\coloneqq k^{-1}$ and let $h_n$ denote the $n$th truncation of $h$. Then
\begin{align*}
\lim_{n\rightarrow\infty}n\norm{h_n-(v_n'h_n) v_n}_{\infty}=1. \quad\text{(Conjecture)}
\end{align*}
The sequence $h_n$ arises in combination with $\mu_n$ in the asymptotic estimate,
\begin{align*}
h_n'\mu_n = \sum_{k=1}^{n} \frac{\mu(k)}{k} = o(1),
\end{align*}
which is equivalent to the prime number theorem~\cite{shapiro1949some}.


\section{Acknowledgments}
We thank the anonymous referees for the extraordinary, insightful comments. One anonymous referee 
contributed most of Section~\ref{subsec:cos}, all of Section~\ref{subsec:bv}, strengthened the main theorem of section~\ref{sec:main}, strengthened the bounds in
Proposition~\ref{prp:lambdanbound}, and
contributed arguments that greatly improved many proofs.
Computing resources for this research were provided by American Family Insurance.

\begin{thebibliography}{9}

\bibitem{apostol1998introduction}
Tom~M. Apostol, {\em {Introduction to Analytic Number Theory}}, Undergrad.
  Texts Math., Springer, 1998.

\bibitem{benito2008recursive}
Manuel Benito and Juan Varona, {Recursive formulas related to the summation of
  the M{\"o}bius function}, {\em Open Math.\ J.} {\bf 1} (2008), 25--34.

\bibitem{deleglise1996computing}
Marc Del{\'e}glise and Jo{\"e}l Rivat, {Computing the summation of the
  M{\"o}bius function}, {\em Exp.\ Math.} {\bf 5} (1996), 291--295.

\bibitem{gegenbauer1885}
Leopold Gegenbauer, Asymptotisch Gesetze der Zahlentheorie, {\em Denkschriften
  der Akad. Wiss. zu Wien} {\bf 49} (1885), 37--80.

\bibitem{lehman1960liouville}
{R. Sherman} Lehman, {On Liouville's function}, {\em Math.\ Comp.} {\bf 14}
  (1960), 311--320.

\bibitem{shapiro1949some}
Harold Shapiro, Some assertions equivalent to the prime number theorem for
  arithmetic progressions, {\em Comm. Pure Appl. Math.} {\bf 2} (1949),
  293--308.

\end{thebibliography}

\bigskip
\hrule
\bigskip

\noindent{\em 2010 Mathematics Subject Classification:} Primary 11A25; Secondary 11K31.\\
\noindent{\em Keywords:} M\"{o}bius function, Mertens function, Dirichlet convolution, $(-1,0,1)$-valued function, Hermitian matrix.

\bigskip
\hrule
\bigskip

\noindent (Concerned with sequences  \seqnum{A002321},  \seqnum{A008683}, \seqnum{A056594}, \seqnum{A127473}, and \seqnum{A327580}.)

\bigskip
\hrule
\bigskip

\vspace*{+.1in}
\noindent
Received  September 21 2019;
revised versions received  September 22 2019; December 29 2019; July 7 2020.
Published in {\it Journal of Integer Sequences}, August 7 2020.

\bigskip
\hrule
\bigskip

\noindent
Return to
\htmladdnormallink{Journal of Integer Sequences home page}{https://cs.uwaterloo.ca/journals/JIS/}.
\vskip .1in


\end{document}


